
\section{Phänomenologische und Mikroskopische Theorie}

\subsection{Typenunterscheidung}
Supraleiter können in zwei Subklassen eingeteilt werden.
Typ I Supraleiter haben zwei voneinander scharf getrennte Phasen, die supraleitende Meißner-Phase für $T< T_C$ und $H<H_C$ und die normale Phase für $T< T_C$ oder $H<H_C$. Dabei gilt für das angelegte kritische Magnetfeld
\begin{equation}
    H_C(T) = H_0 \cdot \left[1 - \frac{T^2}{T_C^2}\right].
\end{equation}
Typ II Supraleiter haben eine dirtte Phase, genannt Shubnikov-Phase, in welcher das Magnetfeld durch quantisierte Flussschläuche (Abrikosov-Vortices) in den Supraleiter eindringt. Diese Phase existiert für $H_{C1} < H < H_{C2}$, wobei $H_{C1}$ das untere und $H_{C2}$ das obere kritische Magnetfeld bezeichnet. Typ II Supraleiter sind für technische Anwendungen besonders interessant, da sie auch bei hohen Magnetfeldern supraleitend bleiben.
\subsection{Meissner-Ochsenfeld-Effekt und London-Theorie}
Die phänomenologische Beschreibung der Supraleitung wurde 1935 von den Brüdern Fritz und Heinz London formuliert. Ausgehend von der Annahme, dass die supraleitenden Elektronen (Dichte $n_s$, Masse $m$, Ladung $-e$) reibungsfrei beschleunigt werden, lautet die Bewegungsgleichung $m \dot{\vec{v}} = -e \vec{E}$. Mit der Suprastromdichte $\vec{j}_s = -e n_s \vec{v}$ folgt die erste London-Gleichung:
\begin{equation}
    \Lambda \frac{\partial \vec{j}_s}{\partial t} = \vec{E}
\end{equation}
mit $\Lambda = \frac{m}{n_s e^2}$ und $n_s$ der Dichte der supraleitenden Ladungsträger.
Durch Anwendung der Rotation und Einsetzen des Induktionsgesetzes $\nabla \times \vec{E} = -\dot{\vec{B}}$ erhält man nach zeitlicher Integration die zweite London-Gleichung, welche den Meissner-Ochsenfeld-Effekt beschreibt:
\begin{equation}
    \nabla \times \Lambda \vec{j}_s = -\frac{1}{c} \vec{B}
\end{equation}
Kombiniert man dies mit dem Ampèreschen Gesetz $\nabla \times \vec{B} = \mu_0 \vec{j}_s$, ergibt sich die Differentialgleichung $\nabla^2 \vec{B} = \frac{1}{\lambda_L^2} \vec{B}$. Die charakteristische Längenskala, auf der ein Magnetfeld exponentiell in den Supraleiter eindringt, ist die London-Eindringtiefe:
\begin{equation}
    \lambda_L = \sqrt{\frac{m c^2}{\mu_0 n_s e^2}}
\end{equation}
mit $\mu_0$ der magnetischen Feldkonstante und $e$ der Elementarladung.
Durch die Temperaturabhängigkeit der supraleitenden Ladungsträgerdichte $n_s$, ergibt sich für die Eindringtiefe die Abhängigkeit
\begin{equation}
    \lambda_L (T) = \frac{\lambda_L(0)}{\sqrt{1 - (T/T_c)^4}}.
\end{equation}

\subsection{Flussquantisierung}
Der supraleitende Zustand lässt sich durch eine makroskopische Wellenfunktion $\Psi(\vec{r}) = \sqrt{n_s(\vec{r})} e^{i\theta(\vec{r})}$ beschreiben. Der quantenmechanische Stromdichteoperator liefert für die Suprastromdichte:
\begin{equation}
    \vec{j}_s = \frac{n_s q}{m} (\hbar \nabla \theta - q \vec{A})
\end{equation}
wobei $q = -2e$ die Ladung eines Cooper-Paares ist. Tief im Inneren eines massiven Supraleiters verschwindet die Stromdichte ($\vec{j}_s = 0$). Integriert man entlang eines geschlossenen Weges tief im Supraleiter, der einen nicht-supraleitenden Bereich (z.B. ein Loch) umschließt, folgt $\hbar \oint \nabla \theta \cdot d\vec{l} = q \oint \vec{A} \cdot d\vec{l}$. Da die Wellenfunktion eindeutig sein muss, gilt $\oint \nabla \theta \cdot d\vec{l} = 2\pi n$ mit $n \in \mathbb{Z}$. Mit dem Stokesschen Satz $\oint \vec{A} \cdot d\vec{l} = \int \vec{B} \cdot d\vec{S} = \Phi$ folgt die Quantisierung des magnetischen Flusses:
\begin{equation}
    \Phi = n \frac{h}{2e} = n \Phi_0
\end{equation}
mit dem magnetischen Flussquant $\Phi_0 \approx \SI{2.068e-15}{\weber}$.

\subsection{BCS-Theorie und die Energielücke}
Nach der BCS-Theorie bilden Elektronen mit entgegengesetztem Impuls und Spin $(\vec{k}\uparrow, -\vec{k}\downarrow)$ aufgrund einer schwachen, durch Phononen vermittelten anziehenden Wechselwirkung gebundene Zustände, Paare, welche sich wie Bosonen verhalten. Die sogenannten Cooper-Paare haben einen Gesamtspin von $S=0$, und kondensieren in einen makroskopischen Grundzustand. 

Um ein Cooper-Paar aufzubrechen, muss die Bindungsenergie $2\Delta$ aufgewendet werden. Das Anregungsspektrum der Quasiteilchen weist daher eine Energielücke der Breite $2\Delta$ um die Fermi-Energie auf. Die BCS-Zustandsdichte der Quasiteilchen divergiert an den Rändern der Energielücke:
\begin{equation}
    D_S(E) = D_N(E_F) \frac{|E|}{\sqrt{E^2 - \Delta^2}} \quad \text{für } |E| > \Delta
\end{equation}
Die Energielücke ist stark temperaturabhängig. Für $T \to 0$ gilt $2\Delta(0) \approx 3.52 k_B T_c$. In der Nähe der kritischen Temperatur $T_c$ (für Niob $T_c \approx \SI{9.2}{\kelvin}$) fällt die Energielücke stark ab:
\begin{equation}
    \Delta(T) \approx 1.74 \Delta(0) \sqrt{1 - \frac{T}{T_c}}
\end{equation}
Aufgrund dieses wurzelförmigen Verhaltens ändert sich $\Delta(T)$ nahe $T_c$ extrem schnell. Um den Verlauf der $\Delta(T)$-Kurve im Experiment präzise zu erfassen, müssen die Messpunkte nahe $T_c$ deutlich dichter gewählt werden als bei tiefen Temperaturen.

\section{Tunnelprozesse und Josephson-Effekte}

\subsection{Quasiteilchen-Tunneln (SIS, NIS, NIN)}
Trennt man zwei Metalle durch eine dünne Isolatorschicht ($d \approx \SI{1}{\nano\meter}$ bis \SI{3}{\nano\meter}), können Elektronen quantenmechanisch tunneln. 
\begin{itemize}
    \item \textbf{NIN (Normalleiter-Isolator-Normalleiter):} Es ergibt sich ein lineares Ohmsches Verhalten $I \propto V$.
    \item \textbf{NIS (Normalleiter-Isolator-Supraleiter):} Bei $T=0$ fließt kein Strom, bis die Spannung $V = \Delta/e$ erreicht, da vorher keine freien Zustände zur Verfügung stehen.
    \item \textbf{SIS (Supraleiter-Isolator-Supraleiter):} Bei $T=0$ setzt der Tunnelstrom erst bei $V = 2\Delta/e$ abrupt ein. Bei endlichen Temperaturen ($T > 0$) existiert ein kleiner thermischer Leckstrom für $V < 2\Delta/e$.
\end{itemize}

\subsection{Die Josephson-Gleichungen}
Zusätzlich zu den Quasiteilchen können auch intakte Cooper-Paare durch die Barriere tunneln. Nach einem Modell von Feynman lässt sich dies durch zwei schwach gekoppelte makroskopische Wellenfunktionen $\Psi_1$ und $\Psi_2$ beschreiben:
\begin{equation}
    i\hbar \frac{\partial \Psi_1}{\partial t} = E_1 \Psi_1 + K \Psi_2 \quad \text{und} \quad i\hbar \frac{\partial \Psi_2}{\partial t} = E_2 \Psi_2 + K \Psi_1
\end{equation}
wobei $K$ die Kopplungskonstante der Barriere ist. Mit dem Ansatz $\Psi_{1,2} = \sqrt{n_{1,2}} e^{i\theta_{1,2}}$ und der Phasendifferenz $\varphi = \theta_2 - \theta_1$ erhält man nach Trennung von Real- und Imaginärteil die beiden Josephson-Gleichungen:
\begin{equation}
    I = I_c \sin(\varphi)
    \label{eq:jos1}
\end{equation}
\begin{equation}
   \frac{\partial \varphi}{\partial t} = \frac{2e}{\hbar} V
   \label{eq:jos2}
\end{equation}
\cref{eq:jos1} beschreibt den DC-Josephson-Effekt: Ein supraleitender Gleichstrom bis zu einem Maximalwert $I_c$ kann bei $V=0$ fließen. \cref{eq:jos2} beschreibt den AC-Josephson-Effekt: Liegt eine endliche Spannung $V$ an, ändert sich die Phase linear in der Zeit ($\varphi(t) = \varphi_0 + \omega_J t$), was zu einem Wechselstrom mit der Frequenz $f_J = \frac{2e}{h}V$ führt ($\frac{2e}{h}=$\SI{483.6}{\giga\hertz\per\milli\volt}).

\paragraph{DC-Josephsoneffekt}
Für diesen Effekt misst man den Tunnelstrom für $V=0$ bei unterschiedlichen Temperaturen. Durch die 2. Josephson-Gleichung muss \(\dot{\delta} = 0\) und somit \(\delta = \text{const.}\). Da die Phasendifferenz nur vom angelegten Strom abhängt, muss durch die 1. Josephson-Gleichung ein nicht verschwindender und zeitlich konstanter Tunnelstrom \(I_S\) fließen.
Über der kritischen Stromstärke \(I_S\) bricht das Cooper-Paar-Tunneln zusammen und es fällt eine Spannung am Kontakt ab. Die typische $I$-$V$-Kennlinie eines Josephson-Kontaktes ist in \cref{fig:iv_dc_josephson} dargestellt.

\begin{figure}
    \centering
    %% \includegraphics{images/dc_jo.png}
    \includegraphics{example-image-a}
    \caption{Typische $I$-$V$-Kennlinie eines Josephson-Kontaktes. Für $I < I_c$ fließt ein supraleitender Gleichstrom ($V=0$). Bei $I > I_c$ bricht das Cooper-Paar-Tunneln zusammen und es fällt eine Spannung ab.}
    \label{fig:iv_dc_josephson}
\end{figure}

\paragraph{AC-Josephsoneffekt} Für diesen Effekt misst man den Tunnelstrom für eine endliche Spannung $V\neq0$. Nach der 2. Josephson-Gleichung gilt
\begin{equation}
    \delta = \frac{2e}{\hbar} U \cdot t + \delta_0 = \omega_J t + \delta_0
\end{equation}
mit der Josephsonfrequenz \(\omega_J = \frac{2e}{\hbar} U\). Somit ändert sich die Phasendifferenz linear in der Zeit und die 1. Josephson-Gleichung liefert einen Wechselstrom mit der Frequenz \(\omega_J\):
\begin{equation}
    I_S(t) = I_C \cdot \sin (\omega_J t + \delta_0)
\end{equation}

\subsection{Das RCSJ-Modell und das Waschbrettpotential}
Reale Josephson-Kontakte weisen neben dem idealen Tunnelstrom auch eine Kapazität $C$ (durch die Plattenkondensator-Geometrie) und einen Normalwiderstand $R$ (durch Quasiteilchen-Tunneln oder externe Shunts) auf. Das Resistively and Capacitively Shunted Junction (RCSJ) Modell beschreibt den Kontakt als Parallelschaltung dieser drei Elemente. Die Kirchhoffsche Knotenregel liefert:
\begin{equation}
    I = I_c \sin(\varphi) + \frac{V}{R} + C \frac{dV}{dt}
\end{equation}
Mit der zweiten Josephson-Gleichung lässt sich dies in eine nichtlineare Differentialgleichung für die Phase $\varphi$ umschreiben:
\begin{equation}
    \frac{\hbar C}{2e} \ddot{\varphi} + \frac{\hbar}{2eR} \dot{\varphi} + I_c \sin(\varphi) = I
\end{equation}
Diese Gleichung ist mathematisch identisch zur Bewegungsgleichung eines Teilchens der Masse $m \propto C$ mit Reibung $\gamma \propto 1/R$ in einem Waschbrettpotential (Tilted Washboard Potential):
\begin{equation}
    U(\varphi) = -E_J \cos(\varphi) - \frac{\hbar I}{2e} \varphi \quad \text{mit} \quad E_J = \frac{\hbar I_c}{2e}
\end{equation}
Die Dynamik wird maßgeblich durch den dimensionslosen Stewart-McCumber-Parameter $\beta_c$ bestimmt:
\begin{equation}
    \beta_c = \frac{2e I_c R^2 C}{\hbar}
\end{equation}
\begin{itemize}
    \item \textbf{Unterdämpfter Fall ($\beta_c \gg 1$):} Die Masse (Kapazität) dominiert. Das Teilchen kann im Potentialtal ruhen ($V=0$) oder über die Maxima rollen ($V \neq 0$). Dies führt zu einer starken Hysterese in der $I$-$V$-Kennlinie. Solche Kontakte sind typischerweise unshunted.
    \item \textbf{Überdämpfter Fall ($\beta_c \ll 1$):} Die Reibung dominiert. Sobald $I < I_c$ fällt, bleibt das Teilchen sofort im nächsten Minimum liegen. Die Hysterese verschwindet. Dies wird experimentell durch einen parallelen Shunt-Widerstand erzwungen. Im Praktikum ist anhand der $I$-$V$-Kurve zu identifizieren, ob der gemessene Kontakt geshuntet ist.
\end{itemize}

%% Hier noch Temperaturabhängigkeit einfügen


\subsection{Magnetfeldabhängigkeit und Fraunhofer-Muster}
Ein Magnetfeld $\vec{B}$ parallel zur Barriere moduliert die Phase räumlich. Aus der Eichinvarianz folgt der Gradient der Phasendifferenz entlang der Kontaktbreite $x$:
\begin{equation}
    \frac{\partial \varphi}{\partial x} = \frac{2e d}{\hbar} B
\end{equation}
wobei $d = 2\lambda_L + t_{ox}$ die magnetische Eindringdicke (inklusive Oxidschicht $t_{ox}$) ist. Integriert man die lokale Stromdichte $J_c \sin(\varphi(x))$ über einen rechteckigen Kontakt der Breite $w$, erhält man für den maximalen kritischen Strom ein Fraunhofer-Beugungsmuster:
\begin{equation}
    I_c(B) = I_c(0) \left| \frac{\sin\left(\pi \frac{\Phi}{\Phi_0}\right)}{\pi \frac{\Phi}{\Phi_0}} \right|
\end{equation}
mit dem magnetischen Fluss $\Phi = B \cdot w \cdot d$ durch den Kontakt. Die Minima treten bei ganzzahligen Vielfachen des Flussquants auf ($\Phi = n\Phi_0$). 

Diese Ableitung setzt voraus, dass das Eigenfeld des Josephson-Stroms vernachlässigbar ist. Dies ist der Fall für "kurze" Kontakte ($w < \lambda_J$). Die Josephson-Eindringtiefe $\lambda_J$ gibt an, wie tief ein Magnetfeld in den Kontakt parallel zur Barriere eindringen kann:
\begin{equation}
    \lambda_J = \sqrt{\frac{\hbar}{2e \mu_0 d J_c}}
\end{equation}
Aus der experimentellen Bestimmung der Minima-Abstände $\Delta B$ lässt sich bei bekannter Kontaktgeometrie die London-Eindringtiefe $\lambda_L$ des Niobs extrahieren.

\subsection{SQUID (Superconducting Quantum Interference Device)}
Schaltet man zwei Josephson-Kontakte in einem supraleitenden Ring parallel, entsteht ein DC-SQUID. Die makroskopischen Wellenfunktionen interferieren in Abhängigkeit des magnetischen Flusses durch den Ring. Der kritische Gesamtstrom moduliert periodisch mit $\Phi_0$, was SQUIDs zu den empfindlichsten Magnetometern der Welt macht: $I_{max} = 2I_c |\cos(\pi \Phi / \Phi_0)|$.
